% conclusion.tex — HapGraph
% Target: 200+ words

We have presented \hapgraph{}, an open-source Python tool for joint admixture
graph inference and timing estimation from phased SNP data.
By combining an F-statistics-based greedy topology search with a Bayesian
IBD-based timing model, \hapgraph{} addresses a gap that has persisted in
population genetic inference: the lack of a unified framework for estimating
who mixed with whom and when, within a single calibrated analysis.

Simulation experiments across three scenarios demonstrate that \hapgraph{}
achieves zero false positive rate on pure trees; greedy topology recall
reaches 85\% for $K=1$ (S2) and 97.5\% for $K=2$ (S3), with precision 85\%
and 100\% respectively.
Under oracle topology, posterior timing estimates are conservative but
well-calibrated (95--97.5\% empirical 95\% CI coverage for $T$ across $K=1$
and $K=2$ scenarios; target: 85\%).
The key limitation identified in benchmarking — severely under-calibrated
uncertainty for admixture proportions (10--20\% empirical 95\% CI coverage)
— is attributable to the delta-method approximation breaking down under high
genetic drift ($N_e = 1{,}000$) and near $\alpha \approx 0.5$; improved
uncertainty quantification for $\alpha$ is a priority for future development.

Application to the 1000 Genomes Project 2022 dataset yielded eight biologically
interpretable admixture events consistent with published population histories,
demonstrating that \hapgraph{}'s topology inference scales to 26 populations
and produces results suitable for downstream biological interpretation.

Future directions include: extending simulation validation to $K \geq 3$;
applying IBD-based timing estimation to real data following hap-ibd
pre-processing; improving alpha credible interval calibration; and direct
quantitative comparison to \textsc{GLOBETROTTER} and \alder{} timing estimates.
\hapgraph{} is designed as a modular Python package, allowing components
(F-statistics, IBD model, MCMC sampler) to be extended or replaced
independently, facilitating community development and adaptation to new
use cases.
